% !TeX spellcheck = it_IT
\chapter{Risultati}\label{chapter:risultati}
Dopo aver definito le relazioni di vicinanza tra sottogeni 
e aver costruito la rete dei rapporti spaziali, si è voluto verificare se i geni classificati 
come \texttt{VICINI} mostrassero effettivamente pattern coerenti di clustering nel
tessuto, ovvero se la loro prossimità spaziale potesse corrispondere 
a domini di espressione condivisi o correlati. Per questo scopo è stato applicato 
un approccio di \textit{clustering spaziale} basato sul'algoritmo \textit{GraphST},
integrato con la modellizzazione statistica offerta dal pacchetto \textit{Mclust} in ambiente R.


\paragraph{Selezione dei geni di interesse}
Come primo passo, dal dataset originale \texttt{AnnData} 
sono stati selezionati esclusivamente i geni appartenenti al cluster 
principale del grafo dei VICINI, ovvero il gruppo centrale emerso dall'analisi 
grafica descritta nella sezione precedente.  

\begin{lstlisting}[style=mypython,
                   label=code:NN,
                   caption={Creazione anndata con solo geni del cluster principale dei VICINI}],
geneIdsToSaveMainHub = [837,810,183,240,
						578,116,193,700,
						792,892,437,815,
						691,553,278,843]

originalGeneIdsToSaveMainHub=[]
for i in geneIdsToSaveMainHub:
    originalGeneIdsToSaveMainHub.append(original_gene_ids[i])

adata = load_data()
adata_main_hub  = adata[:, originalGeneIdsToSaveMainHub].copy()
\end{lstlisting}

\paragraph{Addestramento del modello GraphST}
Il sottoinsieme così ottenuto viene quindi utilizzato per l'addestramento 
di un modello \textit{GraphST}, una metodologia di analisi non supervisionata 
che integra l'informazione spaziale con l'espressione genica per 
individuare cluster coerenti di spot.
Il modello sfrutta una rappresentazione grafica dei dati per apprendere
una codifica latente delle relazioni locali tra geni e posizioni spaziali.
dopodichè facciamo l'addestramento del modello Con GraphST.

\begin{lstlisting}[style=mypython,
                   label=code:NN,
                   caption={Addestramento modello con GraphSt}],
model = GraphST.GraphST(adata_main_hub, device=device)

adata_main_hub = model.train()
\end{lstlisting}
e infine applichiamo la funzione di clustering che sfrutta McLust per trovare i pattern
\begin{lstlisting}[style=mypython,
                   label=code:NN,
                   caption={Addestramento modello con GraphSt}],
adata_main_hub = ad.read_h5ad("main_hub_close_adata")

Mc_Lust(adata_main_hub, n_cluster, refinement=True)
\end{lstlisting}

I risulati ottenuti purtroppo portano alla conclusione che questa metodologia non ha portato 
al ritrovamento di cluster alternativi all'interno del tessuto.